% Reader question: what does an ideal doubling schedule look like?
% Data: one initially informed daemon; eight total. Every informed daemon
% contacts a distinct uninformed daemon in each reliable synchronous round.
% Counts at rounds 0,1,2,3 are 1,2,4,8: fractions 1/8,1/4,1/2,1.
% This favorable schedule illustrates growth, not random-gossip expectation
% or a delivery deadline. Partitions invalidate any finite global deadline.
\begingroup\makeatletter
\ifdefined\pd@typesetpicturebox
  \let\pgfsys@typesetpicturebox\pd@typesetpicturebox
\fi
\begin{tikzpicture}[pd figure,x=20pt,y=42pt]
  \draw[pd hairline] (0,0) -- (3,0);
  \draw[pd guide] (0,1) -- (3,1);
  \draw[pd legible,line width=.9pt] plot coordinates
    {(0,.125) (1,.25) (2,.5) (3,1)};
  \foreach \r/\f in {0/.125,1/.25,2/.5,3/1}
    \node[pd legible datum] at (\r,\f) {};
  \node[pd label,anchor=east] at (0,0) {$0$};
  \node[pd label,anchor=west,inner xsep=2pt] at (3,1) {$1$};
  \node[pd direct label,anchor=south,text width=65pt,align=center]
    at (1.5,1.06) {ideal doubling};
  \node[pd label,anchor=north] at (1.5,0) {rounds 0--3};
\end{tikzpicture}
\makeatother\endgroup
